This lane establishes the \textbf{one-step RG Taylor remainder bounds} in the dyadic seminorm topology.

\paragraph{Target (one-step RG Taylor remainder bounds).}
There exist constants $r_0>0$, $\kappa\in(0,1)$, $A_2\in\mathbb R$, and $C_u,C_K>0$ such that for every dyadic scale $j$
and all $(u,K)$ in the RG ball $\mathbb B_j(r_0)$,
the RG map $\mathcal R_j:(u,K)\mapsto(u',K')$ admits the expansion
\begin{align}
u'&=u+A_2 u^2+\delta_u(u,K),\label{eq:target_u}\\
K'&=\mathcal L_j K+\delta_K(u,K),\label{eq:target_K}
\end{align}
with $\|\mathcal L_j\|_{\mathrm{op}}\le\kappa$ and bounds
\begin{align}
|\delta_u(u,K)|&\le C_u\big(|u|^3+|u|\|K\|_j\big),\label{eq:target_du}\\
\|\delta_K(u,K)\|_{j+1}&\le C_K\big(|u|^2+\|K\|_j^2\big).\label{eq:target_dK}
\end{align}
All constants must be \emph{uniform in $j$}; along the tuned dyadic continuum limit this is the precise form of lattice-spacing uniformity (cf. the chart-truncated setup in Definition~\ref{def:RG_good}).

\paragraph{Strategy.}
We decompose a single RG step into:
\[
\begin{aligned}
\text{(i) fluctuation integration}\ &+\ \text{(ii) reblocking/rescaling}\\
&+\ \text{(iii) extraction of the marginal coordinate}.
\end{aligned}
\]
Each submap is bounded in the polymer seminorm topology, and the Taylor bounds follow from:
\begin{itemize}
\item uniform product estimates in the polymer algebra,
\item uniform cumulant/cluster bounds for local observables under the fluctuation measure,
\item finite-range locality (scale-uniform overlap counting).
\end{itemize}

We work in the standard polymer representation used throughout the project. For completeness we spell out the minimal structure needed.

Fix a dyadic scale $j$. Let $\mathbb B_j$ denote the partition of space into disjoint $j$-blocks (side length $L_j$).
A $j$-\emph{polymer} $X$ is a finite connected union of $j$-blocks.
Let $|X|$ be the number of blocks in $X$.

A \emph{$j$-activity} is a map $K:\mathscr P_j\to\mathcal F$ assigning to each polymer $X$ a local functional $K(X)$
depending only on the gauge field variables in a fixed neighborhood $\mathrm{Nbd}(X)$.

We denote by $\mathscr K_j$ the Banach space of such activities equipped with a scale-$j$ seminorm $\|\cdot\|_j$
that controls:
\begin{itemize}
\item locality (support restricted to $\mathrm{Nbd}(X)$),
\item size (weighted by a factor $A^{|X|}$),
\item and a finite number of field derivatives (to support Taylor expansions).
\end{itemize}

\begin{definition}[Seminorm axioms (multiplicative algebra and locality counting)]\label{ass:seminorm_axioms}
There exist constants $C_{\mathrm{mult}}\ge 1$, $C_{\mathrm{count}}\ge 1$, and $A>1$ (all scale-independent) such that:
\begin{enumerate}[label=(S\arabic*)]
\item (\textbf{Algebra/product bound}) For activities $K_1,K_2\in\mathscr K_j$,
\[
\|K_1\circ K_2\|_j\ \le\ C_{\mathrm{mult}}\|K_1\|_j\,\|K_2\|_j,
\]
where $\circ$ denotes the polymer product/convolution used to assemble disjoint polymers.
\item (\textbf{Local overlap counting}) For any fixed polymer $X$,
the number of polymers $Y$ with $\mathrm{Nbd}(Y)$ overlapping a given neighborhood is bounded by $C_{\mathrm{count}}^{|X|}$.
\item (\textbf{Monotonicity under restriction}) Restricting an activity to a subpolymer does not increase the seminorm.
\end{enumerate}
\end{definition}

\paragraph{Comment.}
These properties are standard for the admissible polymer seminorms used here (polymer weights + derivative control).
We state them explicitly to track constants in subsequent combinatorial sums.

A single RG step integrates out fluctuations at scale $j$ and produces a new effective density at scale $j+1$.

Let $\mu_{j+1}(\cdot\mid U_{\Sigma})$ be the exact Wilson one-step conditional fluctuation law at scale $j+1$, with the block-skeleton field $U_{\Sigma}$ held fixed as in Theorem~\ref{thm:wilson_one_step_factorization}.
By that theorem this conditional law factorizes over $(j+1)$-blocks; after imposing the local good-block truncation the same block-product structure persists.
Define the fluctuation integration operator $\mathbb E_{j+1}$ acting on local functionals by expectation with respect to $\mu_{j+1}(\cdot\mid U_{\Sigma})$, with the dependence on the frozen skeleton data suppressed from the notation.

\begin{proposition}[Uniform finite-range and moment bounds for the exact/chart-truncated one-step law]\label{ass:fluctuation_bounds}
There exist constants $R_{\mathrm{fr}}=0$ and $C_{\mathrm{mom}}\ge 1$ (scale-independent) such that:
\begin{enumerate}[label=(F\arabic*)]
\item (\textbf{Finite range}) Conditional on the fixed block-skeleton field, functionals supported in disjoint $(j+1)$-block neighborhoods are independent under $\mu_{j+1}$. In particular one may take $R_{\mathrm{fr}}=0$ in the counting arguments.  The same statement holds for the chart-truncated law because the good-block event factorizes blockwise.
\item (\textbf{Uniform moment/cumulant bounds}) For every local polynomial functional $P$ of degree $\le p_{\max}$ supported in one $(j+1)$-block neighborhood,
\[
\big|\mathrm{Cum}_n(P,\dots,P)\big|\ \le\ C_{\mathrm{mom}}^{n}\,\|P\|_{\mathrm{loc}}^{n}
\qquad\text{for }n\le n_{\max},
\]
with $\|\cdot\|_{\mathrm{loc}}$ a fixed local norm and $p_{\max},n_{\max}$ fixed constants.
\end{enumerate}
\end{proposition}

\paragraph{Comment.}
These are the uniform-in-scale analytic inputs for the one-step RG analysis.  The finite-range component is the exact Wilson block-factorization statement of Theorem~\ref{thm:wilson_one_step_factorization}, while the moment/cumulant component is proved on the chart-truncated law from Brascamp--Lieb / log-concavity on the small-curvature chart event.  Section~\ref{sec:patching} then identifies the full exact Wilson recursion (Theorem~\ref{thm:wilson_exact_recursion}) and quantifies its difference from the chart-truncated recursion by Theorem~\ref{thm:Wilson_patch_transport}.

We represent the scale-$j$ effective density as
\begin{equation}\label{eq:RG_density_form}
\mathcal Z_j(U)\ =\ e^{-u\,V_j(U)}\ (1+K)(U),
\end{equation}
where $V_j$ is the chosen marginal local functional (normalized so that $u$ corresponds to the GF coupling up to $O(u^2)$),
and $K$ is the irrelevant activity (polymer expansion remainder).

One RG step produces
\[
\mathcal Z_{j+1}\ =\ \mathrm{Rescale}\circ \mathbb E_{j+1}\bigl[\mathcal Z_j\bigr],
\]
then re-expresses $\mathcal Z_{j+1}$ again in the form \eqref{eq:RG_density_form} to define $(u',K')$.
This re-expression is an \emph{exact} coordinate change once the normalization functional $\Pi$ and the reference marginal $V_{j+1}$ are fixed:
\[
u' := \Pi(\mathcal Z_{j+1}),\qquad
K' := e^{u' V_{j+1}}\,\mathcal Z_{j+1}-1.
\]
Thus the ``original'' one-step map $\mathcal R_j$ is the exact blocked Wilson evolution written in the $(u,K)$ coordinates; the chart-truncated map changes only the one-step fluctuation expectation.

\begin{lemma}[Fluctuation integration is bounded in seminorm]\label{lem:E_bounded}
Assume \ref{ass:seminorm_axioms} and \ref{ass:fluctuation_bounds}.
There exists a scale-independent constant $C_E\ge 1$ such that for any $j$-activity $K\in\mathscr K_j$,
the integrated activity $\mathbb E_{j+1}K$ (viewed at scale $j+1$ after reblocking) satisfies
\[
\|\mathbb E_{j+1}K\|_{j+1}\ \le\ C_E\,\|K\|_j.
\]
Moreover, for products,
\[
\|\mathbb E_{j+1}(K_1\circ K_2)\|_{j+1}\ \le\ C_E\,C_{\mathrm{mult}}\|K_1\|_j\|K_2\|_j.
\]
\end{lemma}

\begin{proof}
Finite range implies that only overlaps within an $O(1)$ neighborhood contribute to a given $(j+1)$-polymer; the number of such overlaps is controlled by $C_{\mathrm{count}}$.
Moment bounds control the expectation of each local functional in the derivative-based seminorm.
Collecting factors yields $C_E=C_E(C_{\mathrm{mom}},C_{\mathrm{count}},R_{\mathrm{fr}})$.
\end{proof}

\begin{lemma}[Explicit large--$\varrho$ control of $C_E$ (checkable form)]\label{lem:CE_checkable}
Assume \ref{ass:seminorm_axioms} and \ref{ass:fluctuation_bounds}. Let $\varrho$ denote the uniform convexity parameter appearing in the fluctuation bounds (in particular, in Lemma~\ref{lem:poly_mom}).
Then there exist explicit constants $A_E\ge 1$ and $B_E>0$ depending only on
\[
(p_{\max},n_{\max},n_{\mathrm{loc}},\delta, C_{\mathrm{count}},R_{\mathrm{fr}},C_{\mathrm{mult}})
\]
such that one may choose
\begin{equation}\label{eq:CE_bound_checkable}
C_E\ \le\ A_E\ +\ B_E\,C_{\mathrm{mom}}(\varrho),
\end{equation}
where $C_{\mathrm{mom}}(\varrho)$ is the moment constant in Lemma~\ref{lem:poly_mom}.
In particular, using the $\varrho$--decay built into Lemma~\ref{lem:poly_mom} (polynomial degree $\le p_{\max}$ and $q\le n_{\max}$), one obtains an explicit threshold $\varrho_\ast$ such that
\[
\varrho\ge \varrho_\ast\ \Longrightarrow\ C_E<4.
\]
\end{lemma}

\begin{proof}
By finite range, the integration of a $(j+1)$-local seminorm reduces to finitely many expectations of local polynomials of degree $\le p_{\max}$ in the fluctuation field, with multiplicity controlled by $C_{\mathrm{count}}$ and neighborhood size controlled by $R_{\mathrm{fr}}$.
Each such local expectation is bounded by the uniform $L^q$ moment bound in Lemma~\ref{lem:poly_mom}, with $q$ taken in $\{1,\dots,n_{\max}\}$ as required by the seminorm axioms and the product rule (hence constants depending only on $(p_{\max},n_{\max},n_{\mathrm{loc}},\delta)$).
Collecting the finite overlap/multiplicity factors gives \eqref{eq:CE_bound_checkable} for explicit $A_E,B_E$.
Finally, Lemma~\ref{lem:poly_mom} supplies the explicit decay of $C_{\mathrm{mom}}(\varrho)$ as $\varrho\to\infty$; solving the inequality $A_E+B_E C_{\mathrm{mom}}(\varrho)<4$ gives a checkable $\varrho_\ast$ and hence the stated implication.
\end{proof}

\begin{lemma}[Linearized RG map contracts irrelevants]\label{lem:L_contraction}
Assume \ref{ass:seminorm_axioms} and \ref{ass:fluctuation_bounds}.
There exists $\kappa\in(0,1)$ and scale-independent $C_L\ge 1$ such that for the linearized map $\mathcal L_j$,
\[
\|\mathcal L_j K\|_{j+1}\ \le\ \kappa\,\|K\|_j
\qquad\text{for all }K\in\mathscr K_j,
\]
uniformly in $j$.
\end{lemma}

\paragraph{Mechanism.}
$\kappa$ is the combination of (i) geometric contraction under reblocking/rescaling (irrelevant dimensions) and
(ii) the finite-range smoothing under $\mathbb E_{j+1}$.
This is the precise place where the RG scaling dimension enters.

\begin{lemma}[Irrelevant generation from the marginal interaction is $O(u^2)$]\label{lem:K_from_u}
Assume \ref{ass:seminorm_axioms} and \ref{ass:fluctuation_bounds}.
There exists a scale-independent constant $C_{u\to K}>0$ such that, setting $K\equiv 0$ and performing one RG step,
the resulting irrelevant activity satisfies
\[
\|K'(u,0)\|_{j+1}\ \le\ C_{u\to K}\,u^2
\qquad\text{for all }|u|\le r_0.
\]
\end{lemma}

\begin{proof}
Expand $e^{-uV_j}=1-uV_j+\tfrac12 u^2 V_j^2+R_3$.
The extraction map removes the relevant/marginal components at orders $u$ and $u^2$; the remainder contributes to $K'$.
Uniform cumulant bounds control $\mathbb E_{j+1}[V_j^2]$ in the seminorm and yield the stated $u^2$ bound.
\end{proof}

\begin{lemma}[Marginal coordinate update with explicit cubic remainder]\label{lem:u_flow}
Assume \ref{ass:fluctuation_bounds} and that the extraction map defines $u'$ by matching a fixed normalization functional $\Pi$:
\[
u':=\Pi(\mathcal Z_{j+1}),
\]
where $\Pi$ is linear and scale-independent.
Then there exists a constant $A_2$ (scale-independent) and $C_{u,3}>0$ such that
\[
u'(u,0)=u + A_2 u^2 + r_u(u),
\qquad |r_u(u)|\le C_{u,3}\,|u|^3.
\]
\end{lemma}

\paragraph{Identification of $A_2$.}
$A_2$ is given by a (scale-independent) second-order cumulant of the normalized marginal observable.
In the chosen GF scheme, fixes $A_2=2b_0\log2$ with $b_0>0$.

\begin{lemma}[Mixed remainder bounds]\label{lem:mixed_bounds}
Assume \ref{ass:seminorm_axioms}, \ref{ass:fluctuation_bounds}, and the contraction Lemma~\ref{lem:L_contraction}.
There exist scale-independent constants $C_{uK},C_{K^2}>0$ such that, for $(u,K)\in\mathbb B_j(r_0)$,
\begin{align*}
|u'(u,K)-u'(u,0)|&\le C_{uK}\,|u|\,\|K\|_j,\\
\|K'(u,K)-K'(u,0)-\mathcal L_j K\|_{j+1}&\le C_{K^2}\,\|K\|_j^2.
\end{align*}
\end{lemma}

\begin{proof}
Use the algebra bound \ref{ass:seminorm_axioms}(S1) to control products of $K$ with the $u$-dependent factor,
and Lemma~\ref{lem:E_bounded} to control expectations. The difference from the linearization is a quadratic remainder
bounded by $C_{\mathrm{mult}}C_E\,\|K\|_j^2$.
The $uK$ bound follows because extracting $u'$ depends linearly on the density and $K$ enters multiplied by a factor $e^{-uV_j}=1+O(u)$.
\end{proof}

\begin{theorem}[One-step analytic RG bounds with tracked constants]\label{thm:RG_analytic_bounds}
Assume \ref{ass:seminorm_axioms} and \ref{ass:fluctuation_bounds}, and assume the extraction maps are chosen as in Lemma~\ref{lem:u_flow}.
Define $\mathcal L_j$ as the linearization of $K\mapsto K'$ at $(u,K)=(0,0)$.

Then the RG map admits the expansion \eqref{eq:target_u}--\eqref{eq:target_dK} with:
\[
\kappa\ \text{from Lemma \ref{lem:L_contraction}},\qquad
C_K:=C_{u\to K}+C_{K^2},\qquad
C_u:=C_{u,3}+C_{uK},
\]
and $A_2$ from Lemma~\ref{lem:u_flow}$.$
All constants are scale-independent.
\end{theorem}

\begin{proof}
Write $u'(u,K)=u'(u,0)+[u'(u,K)-u'(u,0)]$ and apply Lemmas \ref{lem:u_flow} and \ref{lem:mixed_bounds}.
Similarly, write
\[
K'(u,K)=\mathcal L_j K + K'(u,0) + \big[K'(u,K)-K'(u,0)-\mathcal L_jK\big],
\]
and apply Lemmas \ref{lem:K_from_u} and \ref{lem:mixed_bounds}, together with Lemma \ref{lem:L_contraction}.
\end{proof}

Uniformity in $j$ (hence $a/L$) is guaranteed if:
\begin{enumerate}[label=(U\arabic*)]
\item the seminorm axioms \ref{ass:seminorm_axioms} hold with constants independent of scale (scale-invariant definition after rescaling);
\item the fluctuation measure bounds \ref{ass:fluctuation_bounds} hold with constants independent of scale (the exact one-step block factorization and BL/moment bounds uniform under the small-curvature chart);
\item the extraction functional $\Pi$ and normalization of $V_j$ are chosen scale-independently (fixed GF scheme and dyadic rescaling).
\end{enumerate}

\paragraph{What is proved in this section.}
The analytic inputs used downstream are proved within this manuscript.
In particular, the fluctuation-measure bounds (Proposition~\ref{ass:fluctuation_bounds}) combine:
(i) the exact Wilson one-step block factorization of Theorem~\ref{thm:wilson_one_step_factorization}, which supplies the finite-range/block-independence input;
(ii) the chart-truncated uniform convexity theorem (Theorem~\ref{thm:uniform_convexity}); and
(iii) the polynomial cumulant bound (Theorem~\ref{thm:poly_cumulant}).
The linear contraction estimate (Lemma~\ref{lem:L_contraction}) is proved in Lemma~\ref{lem:kappa} using the irrelevance scaling gap
(Lemma~\ref{lem:Delta_min}).
All subsequent bounds are algebraic/combinatorial consequences in the dyadic seminorm topology.